\textbf{Proof.}
Fix a bounded potential-outcome array and \(B\in\mathcal B_{\kappa,\overline a}\).  Put
\[
X_B=n^{-1}B^{\top}(D-\mathbb E_\nu D),
\qquad
\Delta_B=-(\widehat\gamma^{\mathrm{HT}}(B)-\gamma_B^*(Y))^{\top}X_B.
\]
Then \(\widehat\tau_B^{\mathrm{CV}}=\widehat\tau_B^{\mathrm{or}}+\Delta_B\).  The exposure-rate condition, \(|Y_i(a)|\le1\), \(|b_{a,i}|\le\overline a\), and \(S_B\succeq n\kappa I_2\) imply, entry by entry,
\[
\left|\frac{Y_i(a)}{\pi_{i,a}^2}
\{\Omega B S_B^{-1}\}_{i,a;l}\right|
\le \overline z\,n^{\alpha+2\beta-1},
\qquad
\overline z=\frac{2\overline a\,\overline d}{\underline\pi^2\kappa}. \tag{1}
\]
The constant is independent of \(B\) and of the outcome array.

For centered indicators having dependency neighborhoods of size at most \(\overline d n^\alpha\), expansion of a fourth power leaves only index quadruples whose dependency graph has no isolated vertex.  There are at most a constant times \(n^2(\overline d n^\alpha)^2\) such quadruples.  Applying this count to the two coefficient coordinates using (1), and to the two coordinates of \(X_B\), yields constants depending only on \(\underline\pi,\overline d,\overline a,\kappa\) such that
\[
\sup_B\mathbb E_\nu
\lVert\widehat\gamma^{\mathrm{HT}}(B)-\gamma_B^*(Y)\rVert_2^4
\le C_1n^{6\alpha+8\beta-2},
\qquad
\sup_B\mathbb E_\nu\lVert X_B\rVert_2^4
\le C_2n^{2\alpha-2}. \tag{2}
\]
Cauchy--Schwarz and (2) give
\[
\sup_B n\mathbb E_\nu\Delta_B^2
\le C_3n^{4\alpha+4\beta-1}\longrightarrow0. \tag{3}
\]

For completeness, the covariance term cannot be discarded merely from (3).  The analogous third-moment dependency count gives, uniformly in \(B\),
\[
\mathbb E_\nu|\widehat\tau_B^{\mathrm{or}}-\tau|^3
\le C_4n^{7\alpha/2+3\beta-2},
\]
\[
\mathbb E_\nu\lVert\widehat\gamma^{\mathrm{HT}}(B)-\gamma_B^*(Y)\rVert_2^3
\le C_5n^{5\alpha+6\beta-2},
\qquad
\mathbb E_\nu\lVert X_B\rVert_2^3
\le C_6n^{2\alpha-2}. \tag{4}
\]
Hölder's inequality applied to (4) yields
\[
\sup_B n\left|
\operatorname{Cov}_\nu(\widehat\tau_B^{\mathrm{or}},\Delta_B\mid Y)
\right|
\le C_7n^{7\alpha/2+3\beta-1}\longrightarrow0, \tag{5}
\]
because \(7\alpha/2+3\beta<4\alpha+4\beta<1\).  Expanding
\[
\operatorname{Var}_\nu(\widehat\tau_B^{\mathrm{or}}+\Delta_B\mid Y)
-\operatorname{Var}_\nu(\widehat\tau_B^{\mathrm{or}}\mid Y)
=\operatorname{Var}_\nu(\Delta_B\mid Y)
+2\operatorname{Cov}_\nu(\widehat\tau_B^{\mathrm{or}},\Delta_B\mid Y)
\]
and using (3)--(5) proves the uniform conditional equivalence.  These are precisely the boundedness, nondegeneracy, positivity, and dependency-neighborhood estimates behind Zhao's published \(4\alpha+4\beta<1\) variance-equivalence regime, now with constants uniform over the explicitly coordinate-bounded class.

The fixed-array oracle variance is obtained by completing the square:
\[
\operatorname{Var}_\nu(\widehat\tau_B^{\mathrm{or}}\mid Y)
=n^{-2}Y^{\top}
\left\{
\Pi^{-1}\Omega\Pi^{-1}
-\Pi^{-1}\Omega BS_B^{-1}B^{\top}\Omega\Pi^{-1}
\right\}Y. \tag{6}
\]
Under ass:iid-potential-pairs, \(\mathbb E_{H^{\otimes n}}(YY^{\top})=Q(c_H)\).  Taking the expectation of (6), applying \(\mathbb E(Y^{\top}MY)=\operatorname{tr}\{M\mathbb E(YY^{\top})\}\), and invoking def:gain-loss gives exactly \(h(c_H)-g_B(c_H)=\ell_B(c_H)\).  The uniform deterministic remainder above remains \(o(n^{-1})\) after expectation over \(H^{\otimes n}\). \(\square\)